\providecommand{\mR}{\mathfrak{R}}
\providecommand{\ma}{\mathfrak{a}}
\providecommand{\mb}{\mathfrak{b}}
\providecommand{\mc}{\mathfrak{c}}
\providecommand{\mm}{\mathfrak{m}}
\providecommand{\mt}{\mathfrak{t}}

\subsection*{\S{}6. Théorie des idéaux sous l'hypothèse de la condition de chaîne des diviseurs}

Soit donné un anneau (commutatif) $\mR$ pour lequel l'axiome I de la condition de chaîne des diviseurs est satisfait; de plus, dans ce qui suit, soit toujours donné un bon ordre fixé de tous les éléments de $\mR$. Par là est en même temps donné aussi un bon ordre de tous les idéaux de $\mR$. En effet, les deux hypothèses entraînent aussitôt que tout idéal de $\mR$ possède une base d'idéal composée d'un nombre fini d'éléments. Si l'on passe donc du bon ordre des éléments de $\mR$ à un ordre quasi-lexicographique de toutes les sous-parties finies de $\mR$,\footnote{Il faut entendre par là ce qui suit: les éléments $a_1,\ldots,a_n$ d'une sous-partie finie $\mathfrak W$ soient désignés de telle sorte que l'ordre des indices concorde avec le bon ordre donné dans $\mR$. Toute sous-partie à $n$ éléments précède une telle sous-partie à $m>n$ éléments. Si deux sous-parties ont même cardinal, la première est dite antérieure lorsque le premier élément où elles diffèrent précède, dans le bon ordre, l'élément correspondant de l'autre. Ainsi un premier élément est assigné à tout système de sous-parties finies. Pour un bon ordre donné de $\mR$, aucun autre postulat du choix n'est donc nécessaire; en particulier, si $\mR$ est dénombrable, comme dans le cas du corps de nombres algébriques, aucun postulat du choix n'est en jeu.} et si l'on assigne à chaque idéal, comme base distinguée, la première, dans le bon ordre des sous-parties finies, parmi les différentes bases possibles, alors les idéaux sont bien ordonnés en même temps que les sous-parties finies.

\textbf{Théorème I.} \emph{Sous l'hypothèse de la condition de chaîne des diviseurs, tout idéal de $\mR$ admet une représentation comme plus petit commun multiple d'un nombre fini d'idéaux irréductibles, c'est-à-dire d'idéaux qui ne peuvent pas se représenter comme plus petit commun multiple de deux diviseurs propres.}

Pour la démonstration, il faut montrer ceci: si le théorème I n'est pas satisfait pour un idéal $\mm$, alors $\mm$ possède un diviseur propre pour lequel le théorème I n'est pas non plus satisfait; on peut en tirer, contrairement à la condition de chaîne des diviseurs supposée, une chaîne de diviseurs qui ne s'interrompt pas après un nombre fini de termes. En effet, $\mm$ doit être réductible, car sinon $\mm=[\mm]$ fournirait la représentation voulue. Si $\mm=[\ma,\mb]$, le théorème I ne peut pas être satisfait simultanément pour $\ma$ et pour $\mb$, car sinon il s'ensuivrait une représentation correspondante pour $\mm$. Il existe donc des diviseurs propres de $\mm$ pour lesquels le théorème I n'est pas satisfait; soit $\ma_1$ le premier d'entre eux dans le bon ordre des idéaux. En construisant de même, à partir de $\ma_1$, un diviseur propre $\ma_2$, et en général, à partir de $\ma_i$, un diviseur propre $\ma_{i+1}$, on arrive à une chaîne de diviseurs bien déterminée, qui ne s'interrompt pas après un nombre fini de termes, contre l'hypothèse.\footnote{Le théorème I vaut évidemment de même pour les domaines de modules, si la condition de chaîne des diviseurs est supposée pour le système de tous les modules. Il a un caractère purement ensembliste; il est en même temps le seul théorème dans la démonstration duquel le bon ordre des idéaux est utilisé.}

À cause de la condition de chaîne des diviseurs supposée, on peut, d'après \S{}5, 2, parler simplement d'idéaux primaires. Le lien entre primaire et irréductible est donné par:

\textbf{Théorème II.} \emph{Sous l'hypothèse de la condition de chaîne des diviseurs, tout idéal irréductible est primaire; en d'autres termes, tout idéal non primaire est réductible.}

Comme, lors du passage à l'anneau quotient $\mR/\mm$, la condition de chaîne des diviseurs est conservée par homomorphie, comme à la représentation de $\mm$ comme plus petit commun multiple correspond la représentation de l'idéal nul, et comme, à cause du premier théorème d'isomorphisme (\S{}4, 3), aux diviseurs primaires de $\mm$ correspondent de nouveau de tels idéaux dans $\mR/\mm$, et réciproquement, on peut se limiter à la décomposition de l'idéal nul dans l'anneau quotient. Puisque celui-ci est un anneau de même généralité, on peut d'emblée supposer que l'idéal à décomposer est l'idéal nul de $\mR$.

Soit donc l'idéal nul de $\mR$ non primaire, de sorte qu'il existe au moins un couple d'idéaux $\ma,\mb$, où $\mb$ peut encore être supposé idéal principal, tel que
\[
\ma\ne(0),\qquad \mb^x\ne(0)\quad\text{pour tout }x,
\qquad \ma\mb=(0).
\]
Formons la chaîne, qui par hypothèse s'interrompt après un nombre fini de termes, des quotients d'idéaux:
\[
\ma,\ \ma:\mb,\ldots,\ma:\mb^\nu,\ldots;
\]
soit par exemple $\mt=\ma:\mb^{m-1}=\ma:\mb^m=\cdots$ égal à tous les idéaux suivants de la chaîne; donc $\mt=\mt:\mb$, ou bien $\mb$ est premier à $\mt$. De plus, $\mt$, comme diviseur de $\ma$, est différent de l'idéal nul, et il en va de même de $\mb^x$ pour tout exposant $x$. Pour démontrer le théorème II, il suffit donc d'établir une représentation
\[
(0)=[\mt,\mb^{m+1}].
\]
Par définition,
\[
\mb^{m-1}\mt\equiv0\pmod{\ma},\qquad\text{donc}\qquad \mb^m\mt=(0),
\]
il reste donc seulement à montrer:
\[
\mc=[\mt,\mb^{m+1}]\equiv0\pmod{\mb^m\mt}.
\]
Cela résulte des hypothèses que $\mb$ est un idéal principal et qu'il est premier à $\mt$. Tout élément $c$ de $\mc$ admet, à cause de la divisibilité par $\mb^{m+1}$, une représentation, où $n$ est compris comme symbole d'un entier,
\[
 c=k b^{m+1}+n b^{m+1}=r b^m,
\]
où $r$ désigne maintenant de nouveau un élément de $\mR$. De $c\equiv0\pmod{\mt}$ il résulte donc $r b^m\equiv0\pmod{\mt}$, et par suite $r\equiv0\pmod{\mt}$, puisque $\mt:\mb=\mt$. Mais alors $c\equiv0\pmod{\mt\mb^m}$, et le théorème II est démontré.

\textbf{Théorème III.} \emph{Sous l'hypothèse de la condition de chaîne des diviseurs, tout idéal admet une représentation la plus courte comme plus petit commun multiple d'un nombre fini de composantes primaires les plus grandes, appartenant à des idéaux premiers différents.}

Pour obtenir une telle représentation, il suffit de remplacer, chaque fois que c'est possible, la représentation par idéaux irréductibles qui existe d'après le théorème I, donc par idéaux primaires d'après le théorème II, par une représentation la plus courte. Si l'on rassemble les idéaux primaires appartenant aux mêmes idéaux premiers, on obtient d'après \S{}5, 3 la représentation cherchée. Les composantes primaires peuvent être qualifiées de plus grandes, car le plus petit commun multiple de quelconques d'entre elles n'est plus primaire.\footnote{Les idéaux premiers associés et les composantes isolées sont alors déterminés de façon univoque. Voir \emph{Idealtheorie} ou W. Krull, Math. Ann. 90 (1923), p. 55--64.}
